\section{The resulting configuration}\label{sec:config}

This section is the practical answer for this machine, chosen in the priority order accuracy, then
usable context, then speed. It selects a complete serving configuration, not the better weight quantization, and it is a decision for one host and one workload (a single-user coding
agent). The rest of the report does not depend on it.

\begin{table}[t]
\centering\small
\caption{The configuration that now serves on this host, and its tested fallback. Both hold the native
262{,}144-token window in one slot.}\label{tab:config}
\begin{tabularx}{\linewidth}{@{}lXX@{}}
\toprule
 & \textbf{serving} & fallback \\
\midrule
weights & \Qfour{} (16.4\,GiB) & \Qsix{} (20.5\,GiB) \\
KV cache & \flag{q8\_0} & \flag{q4\_0} \\
drafter & DFlash2, draft depth 7, drafter on GPU1 & built-in MTP, depth 4 \\
engine & upstream llama.cpp b10975 (E-C) & MTP build (E-A) \\
split & layer, \flag{-ts 58,42} & layer, \flag{-ts 58,42} \\
decode at $\sim$100K input, thinking & 29.8\tps{} [23.8, 51.3] & 16.3\tps{} [16.1, 16.6] \\
decode at 246K filled, greedy code (\flag{q4\_0} KV probe) & 27.7\tps & 17.3\tps \\
AA-LCR / GPQA Diamond & 87/100 \,/\, 175/198 & 22/30 \,/\, 48/50 (subsets) \\
peak host RAM, 252K-token soak & 1.7\,GiB & 5.4\,GiB \\
\bottomrule
\end{tabularx}
\end{table}

\begin{verbatim}
llama-server -m Qwen3.8-27B-UD-Q4_K_XL.gguf -ngl 99 -sm layer -ts 58,42 \
  -c 262144 -fit off -fa on -ctk q8_0 -ctv q8_0 -b 2048 -ub 512 \
  -np 1 -ctxcp 4 -cram 0 \
  --spec-type draft-dflash -md Qwen3.8-27B-DFlash2-Q4_K_M.gguf \
  -ngld 99 --spec-draft-n-max 7 -devd CUDA1
\end{verbatim}

\paragraph{How the decision went.} After the first battery the choice was \Qsix{} with MTP and a
4-bit cache: the highest-fidelity build that reaches the full window, since speed did not separate the
builds and divergence did. The second battery changed the picture on three counts. No accuracy
instrument separated \Qsix{} with a 4-bit cache from \Qfour{} with an 8-bit cache at the full window
(\cref{sec:aalcr}), and outside evidence makes the 4-bit cache the riskier unknown at depth
(\cref{sec:kvq}). DFlash2 was 45--65\,\% faster at every depth (\cref{sec:s15speed}). And it used 3.6\,GiB
less host RAM on a 14\,GiB machine whose only outage had been host RAM (\cref{sec:hostram}). On
divergence \Qfour{} is measurably further from the reference than \Qsix; whether a user would ever
notice is exactly what the task instruments could not show. I weighed a measured speed and memory gain
against an unmeasurable accuracy difference, and kept the 6-bit line one restart away.

\paragraph{What did not survive first contact.} The first cut-over failed. The serving image (E-A)
lists \flag{draft-dflash} in its help text and cannot parse the drafter file: the same ``expected 81,
got 58'' signature that had produced a silent zero a month earlier. The winner needed a different
engine build, not a different flag. The second attempt, on build E-C, passed the gate that the outage
had created: a 44-turn soak to 252{,}525 tokens under serving sampling, peak host memory 1{,}714\,MiB,
no swap growth, rollbacks reprocessing 4 tokens, and recovery from a forced kill to a healthy server
at the full window.

\paragraph{If your hardware differs.} The numbers will not transfer. The procedure should: (1)~compute
the KV budget from the architecture; (2)~sweep the layer split for \emph{your} build, cache precision
and drafter, keeping the fastest split that loads; (3)~accept a window only after a deep prefill
\emph{and} a generation; (4)~prefer the gentlest cache precision that holds your window; (5)~run the
drafter at its design depth and measure at your real context depth with greedy fixed prompts;
(6)~soak-test a growing multi-turn conversation while watching host RAM before trusting any of it.
